% Standalone source: latexmk -pdf writeup.tex
% All proofs, tables, and bibliography entries are contained in this file.
\documentclass[11pt]{article}
\usepackage[a4paper,margin=27mm]{geometry}
\usepackage[T1]{fontenc}
\usepackage{lmodern,microtype}
\usepackage{amsmath,amssymb,amsthm,mathtools}
\usepackage{booktabs,tabularx,array,enumitem}
\usepackage{xurl}
\usepackage[colorlinks=true,linkcolor=blue,citecolor=blue,urlcolor=blue]{hyperref}
\hypersetup{pdftitle={The dimension dependence of local grid coloring},pdfauthor={}}
\newtheorem{theorem}{Theorem}[section]
\newtheorem{lemma}[theorem]{Lemma}
\newtheorem{corollary}[theorem]{Corollary}
\newtheorem{proposition}[theorem]{Proposition}
\theoremstyle{definition}\newtheorem{definition}[theorem]{Definition}
\theoremstyle{remark}\newtheorem{remark}[theorem]{Remark}
\newcommand{\Z}{\mathbb Z}
\newcommand{\R}{\mathbb R}
\newcommand{\E}{\mathbb E}
\newcommand{\Prob}{\mathbb P}
\newcommand{\dist}{\operatorname{dist}}
\newcommand{\LOCAL}{\mathsf{LOCAL}}
\newcommand{\NS}{\mathsf{NS\text{-}LOCAL}}
\newcommand{\BD}{\mathsf{BD\text{-}LOCAL}}
\newcommand{\ls}{\log^*}
\newcommand{\norm}[1]{\lVert#1\rVert}
\newcommand{\abs}[1]{\lvert#1\rvert}
\setlength{\emergencystretch}{2em}
\setlist{itemsep=3pt,topsep=5pt}
\title{The dimension dependence of local grid coloring}
\author{}
\date{30 September 2026}
\begin{document}
\maketitle

\begin{abstract}
We study proper vertex coloring of full $d$-dimensional boxes and standard
tori, treating the dimension as a variable. For every fixed palette
$c\ge4$, a linear dependence on $d$ is necessary even for non-signaling
outcomes with constant global success probability, and hence also for
distributed quantum algorithms with input-independent pre-shared
entanglement. Conversely, four colors suffice in
$2^{O(d)}(1+\log^*n)$ deterministic communication rounds, and in
$2^{O(d)}$ non-signaling locality with exact success or quantum rounds
with high probability. The upper bounds require no supplied coordinate
directions. Classical deterministic and private-randomness algorithms
also have the usual $\Omega_c(\log^*n)$ lower bound, with constants
uniform in dimension.

The lower bound uses a random cyclic Cayley graph with high chromatic
number and induced neighborhoods identical to grid neighborhoods.
The upper bound combines nets at successively doubled scales with
separated box boundaries. We prove the graph construction, the
constant-success non-signaling transfer, and the geometric algorithm
in full. Published non-signaling and quantum symmetry-breaking
primitives are stated explicitly at their use. The remaining gap is
linear versus single-exponential dependence on dimension.
\end{abstract}

{\small\tableofcontents}
\clearpage

\section{Results and conventions}\label{sec:overview}

For integers $d\ge1$ and $L\ge1$, the inputs are the full box
$\{0,\ldots,L-1\}^d$ and the standard torus $(\Z/L\Z)^d$, with
nearest-neighbor edges. Graphs are simple, including for small $L$.
Write $n=L^d$. The maximum degree is at most $2d$ and, when $L\ge3$,
the torus is $2d$-regular. A proper $c$-coloring assigns a member of
$[c]=\{1,\ldots,c\}$ to every vertex and different colors to adjacent
vertices. The palette $c$ is fixed independently of $d$ and $n$ in
the asymptotic comparisons below.

\begin{theorem}[Dimension-dependent bounds]\label{thm:main}
Fix $c\ge4$. On the above grids with
\begin{equation}\label{eq:large}
                         L\ge\max\{128,d^2\},
\end{equation}
the bounds in Table~\ref{tab:main} hold, uniformly in $d$ and $L$.
All upper bounds use only four colors and hold for every $L\ge1$.
The dimension lower bounds hold even when globally consistent signed
coordinate directions are supplied; the upper bounds also work without
these directions. The lower bounds apply to boxes, standard tori, and
tori promised to have even side length.
\end{theorem}

\begin{table}[htbp]
\centering\small
\begin{tabularx}{\linewidth}{@{}lXXl@{}}
\toprule
Model & Lower bound & Upper bound & Upper success\\
\midrule
Deterministic $\LOCAL$
 & $\Omega_c(d+\ls n)$ & $2^{O(d)}(1+\ls n)$ & Certain\\
Private-randomness $\LOCAL$
 & $\Omega_c(d+\ls n)$ & $2^{O(d)}(1+\ls n)$ & Certain\\
Quantum $\LOCAL$
 & $\Omega_c(d)$ & $2^{O(d)}$ & $1-n^{-a}$\\
Non-signaling $\LOCAL$
 & $\Omega_c(d)$ & $2^{O(d)}$ & Certain\\
Bounded-dependence $\LOCAL$
 & $\Omega_c(d)$ & $2^{O(d)}$ & Certain\\
\bottomrule
\end{tabularx}
\caption{Bounds for a fixed palette $c\ge4$. Randomized-model lower
bounds require global success at least $2/3$; deterministic lower
bounds have a harmless additive constant. The quantum upper bound
holds for each fixed requested $a>0$. The private-randomness upper
bound is the deterministic algorithm.}
\label{tab:main}
\end{table}

Here and throughout, success concerns the \emph{entire} coloring, not
just a fixed vertex or edge. The dimension lower bound also holds for
any fixed positive global success probability, and allows shared
randomness independent of the input topology. The additional
randomized $\ls n$ lower bound is stated for independent private
randomness. The factor multiplying $\ls n$ in the deterministic
upper bound depends exponentially on $d$; an additive
$2^{O(d)}+O(\ls n)$ bound is not proved here.

\subsection{What the quantitative arguments add}

Four-coloring grids in $O_d(\ls n)$ deterministic rounds is known
from Brandt et al.~\cite{grid}; the underlying separated-box method
is closely related to Holroyd--Schramm--Wilson~\cite{HSW}.
Our upper-bound argument keeps the dependence on $d$ explicit.
Its key step is to construct a covering net using a sequence of
auxiliary graphs of degree $2^{O(d)}$, instead of asking for a maximal
independent set directly in a grid power of degree $2^{\Theta(d^2)}$.
The separated-box conversion is then reproduced and proved.

Random-Cayley spectral methods are classical; see
Alon~\cite{Alon}. The lower-bound construction combines their
chromatic obstruction with the exclusion of short additive relations,
so that every small induced neighborhood is a signed grid ball.
The transfer to non-signaling outcomes uses the cheating-graph method
of Coiteux-Roy et al.~\cite{NScolor}. We give the complete transfer,
including the amplification to constant global success and the finite
size bookkeeping. The size-dependent classical lower bounds are due
to Linial~\cite{Linial} and Naor~\cite{Naor}; the fast initialization
used for the stronger-model upper bounds is due to
Akbari et al.~\cite{OnlineQuantum} and
Flin--Li--Suomela~\cite{QuantumBreak}.

The two main quantities to compare are thus $\Omega_c(d)$ and
$2^{O(d)}$. The results concern constant palettes and full grid
geometry. Rectangles with arbitrary aspect ratios, grids with holes,
and arbitrary subgraphs are outside the geometric upper-bound theorem.

\subsection{Inputs, distances, and computation}

Vertices know $d,L,n$ and whether the promised input is a box or a
torus. They have distinct identifiers in $[n^C]$ for a fixed known
constant $C\ge1$; all assignments are permitted. Lower bounds use
only identifiers in $[n]$. Local computation and message sizes are
unrestricted. A communication round traverses one nearest-neighbor
edge. Write $\dist_G$ for this graph distance and $\dist_\infty$ for
maximum coordinate distance, using cyclic coordinate distances on a
torus. In these graphs,
\[
              \dist_G(u,v)\le d\,\dist_\infty(u,v).
\]
Thus an $\ell_\infty$ radius $r$ can require $dr$ communication
rounds. Constants in $O(\cdot)$ are independent of $d,n$ unless a
subscript indicates otherwise. The function $\ls$ uses base two and
stops on reaching a value at most one; $\log$ in explicit analytic
inequalities is the natural logarithm.

A \emph{frame} labels the incident directions consistently by
$+e_i,-e_i$, $i\in[d]$. It supplies no absolute coordinates, origin,
or parity bit. Lower bounds remain valid with this extra input.
Absolute coordinates would instead make boxes and even tori
zero-round two-colorable by coordinate-sum parity.

We use the usual radius convention: a vertex's $T$-round information
is its rooted radius-$T$ neighborhood with identifiers, degrees, ports,
and visible input labels. Boundary stubs and their labels are included;
data beyond them are not. Initially knowing neighbor identifiers
changes the bounds by an additive constant. All constructions below
leave an extra layer where a view must be embedded.

\subsection{Non-signaling and quantum conventions}\label{sec:models}

Fix a named processor set of size $n$ and the public parameters.
A \emph{total outcome} assigns a probability distribution on
simultaneous classical vertex outputs to every graph on this set,
including graphs violating the grid promise. It has \emph{non-signaling
locality $T$} if, whenever two inputs have identical joint radius-$T$
views of a set $S$, their output marginals on $S$ agree. Joint views
retain the actual processor names, identifiers, labels, and edges;
they do not merely record separate isomorphism types. Correctness is
required only on promised grid inputs. This is the total-outcome
convention of~\cite{NScolor}.

Deterministic and randomized $T$-round algorithms induce such outcomes.
So do $T$-round quantum algorithms: tracing out all systems outside
the backward communication cones of $S$ leaves a marginal determined
by the operations and initial reduced state in those cones. Operations
outside a cone cannot change that marginal. For the lower bounds, the
initial state may be arbitrarily entangled and may depend on the
processor names and public parameters, but must be independent of
the unknown input topology. The processor set and state are kept
fixed in the comparisons. The quantum upper bounds need no
pre-shared entanglement; local register dimensions are unrestricted.

A bounded-dependence outcome, in addition to this cross-input locality,
has independent output restrictions on sets whose graph distance
exceeds a specified dependence range $k$. Its locality parameter
here bounds both the non-signaling radius and $k$, up to absolute
constant factors. This is a family of finite-input distributions.
A finitely dependent process given only on $\Z^d$ need not extend
to such a family, so the lower bounds below make no assertion about
the dependence range of a standalone infinite-lattice process.

\section{The linear dimension lower bound}\label{sec:lower}

We first construct a graph that looks like a grid throughout the
region visible to an algorithm, yet cannot be colored with its palette.

\begin{lemma}[A quantitative grid-like obstruction]\label{lem:cayley}
Let $d\ge1$, $q\ge2$, and $T\ge0$ be integers. Put
\[
 s=2T+3,\qquad K=\binom{2d+s}{s}.
\]
If
\begin{equation}\label{eq:criterion}
                  \log(64K)<\frac{d}{2(q-1)^2},
\end{equation}
there are a prime $P$ with $8K<P<16K$ and a simple $2d$-regular
graph $H$ on $P$ vertices with $\chi(H)>q$. Every induced closed
radius-$(T+1)$ ball of $H$ is isomorphic to the corresponding ball
in $\Z^d$, preserving the individual signed directions.
\end{lemma}

\begin{proof}
Encode an integer vector by its positive and negative parts.
Stars and bars then gives
\begin{equation}\label{eq:ballcount}
 \#\{z\in\Z^d:\norm{z}_1\le s\}
 \le\#\{w\in\Z_{\ge0}^{2d}:\textstyle\sum_i w_i\le s\}
 =K.
\end{equation}
Bertrand's postulate supplies a prime $P$ strictly between $8K$ and
$16K$. In particular $P>s$, since $K\ge s+1$.
Choose $a_1,\ldots,a_d$ independently and uniformly from $\Z/P\Z$.
For every nonzero $z$ with $\norm{z}_1\le s$, some coefficient $z_i$
is nonzero modulo $P$. Conditioning on all other generators shows
\[
        \Prob\left[\sum_i z_i a_i=0\right]=\frac1P.
\]
The probability that any such relation holds is at most
$(K-1)/P<1/8$. Hence, except on an event of probability less than
$1/8$, we have
\begin{equation}\label{eq:norelation}
       \sum_i z_i a_i\ne0\quad\text{for }0<\norm{z}_1\le s.
\end{equation}

The operator that sums all shifts $\pm a_i$ has Fourier eigenvalues
\[
 \lambda_k=2\sum_{i=1}^d\cos(2\pi k a_i/P),\quad k\in\Z/P\Z;
 \qquad \lambda_0=2d.
\]
Indeed the characters $x\mapsto\exp(2\pi i kx/P)$ form an orthogonal
basis and are eigenvectors of each shift. For fixed $k\ne0$, the
summands in $\lambda_k$ are independent, have mean zero, and lie in
$[-2,2]$. Here is the needed concentration estimate. If $X$ has
these bounds and mean zero, convexity on $[-2,2]$ gives
\[
 \E e^{-tX}\le\cosh(2t)\le e^{2t^2}\qquad(t\ge0).
\]
The second inequality follows by comparing the power series, using
$(2j)!\ge2^j j!$. For a sum of $d$ such independent variables,
Markov's inequality and $t=u/(4d)$, for $u\ge0$, therefore give
\[
 \Prob[\lambda_k\le-u]\le
       \inf_{t\ge0}e^{-tu+2dt^2}=e^{-u^2/(8d)}.
\]
Taking $u=2d/(q-1)$ and summing over the nonzero frequencies bounds
the probability of any violation of
\begin{equation}\label{eq:spectral}
                      \lambda_k>-\frac{2d}{q-1}\quad(k\ne0)
\end{equation}
by
\[
 (P-1)e^{-d/(2(q-1)^2)}<\frac{16K}{64K}=\frac14.
\]
There is consequently a choice of generators satisfying both
\eqref{eq:norelation} and \eqref{eq:spectral}.

For this choice let
\[
 H=\operatorname{Cay}(\Z/P\Z,\{\pm a_1,\ldots,\pm a_d\}).
\]
Since $s\ge3$, the excluded relations include $a_i=0$,
$2a_i=0$, and $a_i=\pm a_j$ for $i\ne j$. Thus the signed
generators are nonzero and distinct: $H$ is simple and $2d$-regular,
and the operator above is its adjacency matrix $A$.

We spell out the chromatic implication. Let $S$ be a nonempty
independent set of density $\alpha=|S|/P$. Since the graph has
positive degree, $0<\alpha<1$. Set $x=1_S-\alpha\mathbf1$.
Then $x\perp\mathbf1$,
\[
 \norm{x}^2=P\alpha(1-\alpha),\qquad
 x^{\mathsf T}Ax=-2dP\alpha^2.
\]
Expanding in the nonconstant Fourier eigenvectors and
using~\eqref{eq:spectral} yields
\[
 -2dP\alpha^2>
 -\frac{2d}{q-1}P\alpha(1-\alpha),
 \qquad\text{hence }\alpha<\frac1q.
\]
Every color class is independent, so a $q$-coloring is impossible.

Finally define $\phi(z)=\sum_i z_i a_i$. Every walk in $H$ lifts to
a sequence of signed coordinate steps, so $\phi$ maps the lattice
radius-$(T+1)$ ball onto the corresponding ball at zero in $H$.
Two distinct points of that lattice ball cannot have the same image:
their difference has norm at most $2T+2$. If their images are
adjacent, subtracting the corresponding signed basis vector gives
a relation of norm at most $2T+3$. By~\eqref{eq:norelation}, this
relation must be the zero vector, so that edge is exactly the
corresponding signed grid edge. This proves both induced adjacency
and preservation of each direction. Translation proves the assertion
at every vertex of $H$.
\end{proof}

\subsection{Transfer without independence of the outputs}

\begin{lemma}[Constant-success amplification]\label{lem:amplify}
Suppose $H$ has $h$ vertices, $\chi(H)>q$, and all its induced
radius-$(T+1)$ balls are signed grid balls. Let $N\ge1$ and suppose
\begin{equation}\label{eq:packing}
 Nh\le L^d,\qquad
 N\le\left\lfloor\frac{L}{2T+5}\right\rfloor^d.
\end{equation}
For every total non-signaling outcome of locality $T$, some legal
named, ID-labeled grid with $n=L^d$ vertices has global $q$-coloring
success probability at most
\[
                         (1-1/h)^N\le e^{-N/h}.
\]
This holds for boxes and standard tori, including even tori when
the prescribed $L$ is even.
\end{lemma}

\begin{proof}
On the fixed named set $[n]$, form $N$ disjoint copies of $H$ and
$n-Nh$ isolated vertices. Keep the same public values of $d,L,n$,
grid-type promise, and fixed distinct identifiers as in the legitimate
inputs. The outcome is defined on this graph even though correctness
is not required there. Replace invalid output symbols by a fixed
member of $[q]$; this preserves locality and cannot turn a successful
legal coloring into a failure.

Call a vertex star bad if its center has the same color as one of
its neighbors. Every realized output vector has at least one bad
star in each copy of $H$. Independently of this vector, an analyst
chooses a uniformly random center in each copy, with independent
choices between copies. Conditional on any fixed vector, the chance
that every selected star is good is at most $(1-1/h)^N$. Averaging
over the vector and then over the choices implies that some
\emph{fixed} selection of $N$ centers has joint probability at most
$(1-1/h)^N$ of all its star checks passing. No independence of
algorithmic outputs has been used.

Let $S$ be the union of these selected stars. In each copy of $H$,
the radius-$T$ neighborhood of its selected star is the
radius-$(T+1)$ ball at the center. Put $b=2T+5$ and partition a
legal grid into $\lfloor L/b\rfloor^d$ disjoint coordinate blocks
of side $b$, ignoring leftover strips. Embed one selected ball
into each of $N$ blocks, centered at coordinate $T+2$ within the
block. Its coordinate support lies between $1$ and $b-2$.
Consequently these embedded balls have no edges between different
blocks and lie away from a box boundary or a torus seam. They have
the full degrees and signed ports needed for the views.

Preserve the actual names and identifiers in each embedded ball,
and assign the unused names and identifiers bijectively to the
remaining grid positions. There are enough vertices by
\eqref{eq:packing}. The joint radius-$T$ view of $S$, including
boundary degrees and ports, is identical on the two inputs.
Non-signaling locality therefore equates their output marginals
on $S$. A globally proper coloring of the genuine grid must pass
all the selected star checks, proving the claim.
\end{proof}

\begin{theorem}[Uniform dimension lower bound]\label{thm:dimension}
Let $q\ge2$ and assume~\eqref{eq:large}. If a total non-signaling
outcome of locality $T$ properly $q$-colors every legal grid with
global probability at least $2/3$, then
\begin{equation}\label{eq:uniformlower}
                 T+2>\frac{d}{512(q-1)^2\log(q+1)}.
\end{equation}
For four colors one also has
\begin{equation}\label{eq:fourlower}
                       T+2>d/256\qquad(d\ge1024).
\end{equation}
Both statements permit a supplied frame and a bipartiteness promise.
\end{theorem}

\begin{proof}
Put $a=(q-1)^2$ and $b=\log(q+1)$, and suppose
$T+2\le d/(512ab)$. Since $T\ge0$, necessarily $d\ge1024ab$.
Also $s=2T+3<d/(256ab)\le d$. The elementary binomial estimate
$\binom{m}{s}\le(em/s)^s$ gives
\[
 \log K\le s\log\frac{e(2d+s)}s
             \le s\log\frac{3ed}s.
\]
The function $x\log(3e/x)$ increases for $0<x\le1$; hence
\[
 \log K\le\frac{d}{256ab}\log(768eab)
             \le\frac{10d}{256a}.
\]
For the last step, use $b\ge\log3$, $\log(768e)\le7b$,
$\log a\le2b$, and $\log b\le b$.
Moreover $\log64\le d/(256a)$ because $d/(256a)\ge4b>\log64$.
Thus
\[
                    \log(64K)\le\frac{11d}{256a}<\frac{d}{2a},
\]
and Lemma~\ref{lem:cayley} applies. As $s\le d$,
\[
                     K\le\binom{3d}{d}\le2^{3d}=8^d,
                     \qquad h=P<16\cdot8^d.
\]
Take $N=2h$. In this nonvacuous regime $d>1024$ and
$2T+5\le d$. Consequently
\[
 Nh=2h^2<512\cdot64^d\le d^{2d}\le L^d,
 \qquad
 N<32\cdot8^d\le d^d\le
       \left\lfloor\frac{L}{2T+5}\right\rfloor^d.
\]
For example, the first elementary inequality follows from
$(d^2/64)^d\ge16^d>512$ for $d\ge1024$; the second follows from
$(d/8)^d\ge128^d>32$. Lemma~\ref{lem:amplify} now bounds success
by $e^{-2}<2/3$, a contradiction.

For~\eqref{eq:fourlower}, suppose $T+2\le d/256$ and $d\ge1024$.
Then $s<d/128$, and monotonicity of
$x\log(e(2d+x)/x)$ gives
\[
 \log(64K)\le\log64+\frac{d}{128}\log(257e)<\frac d{18}.
\]
To check the last strict inequality without numerical approximation,
use $\log64<104/25$ and $\log(257e)<131/20$. The corresponding
margin is exactly
\[
 \frac1{18}-\frac{104/25}{1024}-\frac{131/20}{128}
                         =\frac{37}{115200}>0.
\]
The logarithm bounds follow from the finite rational inequalities
$\sum_{j=0}^{12}(104/25)^j/j!>64$ and
$\sum_{j=0}^{14}(111/20)^j/j!>257$, and the exponential series.
This is~\eqref{eq:criterion} for $q=4$. The same estimates
$s\le d$, $2T+5\le d$, and the same packing with $N=2h$ apply.
\end{proof}

For any fixed $\rho\in(0,1]$, replacing $N=2h$ by
$N=\lceil h\log(2/\rho)\rceil$ makes the success upper bound at most
$\rho/2$. The constant factor depending on $\rho$ is absorbed by
the packing inequalities for sufficiently large $d$. This proves
$\Omega_q(d)$ for every fixed positive success probability in all
the finite-input models of Table~\ref{tab:main}. The proof still
uses only ordinary chromatic number: every run outputs one
simultaneous classical color at every vertex.

\section{Four colors with single-exponential dimension dependence}
\label{sec:upper}

We first describe a deterministic algorithm with identifiers and a
supplied frame. Later we remove the frame and replace the identifiers
used for symmetry breaking by a locally proper initial coloring.

The only general-graph subroutine needed is the following elementary
form of color reduction. A proof, including a uniform finite-table
construction, is in Appendix~\ref{app:reduction}.

\begin{lemma}[Color reduction and greedy independence]\label{lem:reduction}
On a graph of maximum degree at most a known $D\ge1$, a supplied
proper coloring with labels in $[m]$ can be reduced deterministically
to at most $4096(D+1)^6$ colors in $O(1+\ls m)$ rounds.
A maximal independent set can consequently be computed in
$O((D+1)^6+1+\ls m)$ rounds. The constants are absolute;
identifiers are unnecessary once the initial proper coloring is given.
\end{lemma}

The polynomial exponent in this supporting lemma is unimportant for
our dimension bound. Its useful features are uniform constants and
the explicit dependence on the initial palette.

\subsection{Sparse centers from successive scales}

A set of grid vertices is \emph{$a$-separated} if any two distinct
members have maximum-norm distance greater than $a$. It
\emph{$R$-dominates} if every vertex has maximum-norm distance at
most $R$ from the set. All boxes in this section are maximum-norm
balls intersected with the vertex set of the input graph.

\begin{lemma}[Packing]\label{lem:packingupper}
If $S$ is $a$-separated, $a>0$, then any maximum-norm ball of radius
$R$ contains at most $(1+2R/a)^d$ points of $S$. For a torus we
assume $2R+a<L$.
\end{lemma}

\begin{proof}
Put a continuous cube of side $a$ around each selected point.
Their interiors are disjoint, and all these cubes lie in a cube of
side $2R+a$. The volume ratio gives the bound. For a box, the
continuous cubes may extend past the input boundary, which does
not affect this comparison. For a torus, lift the radius-$R$ ball
to coordinate intervals in $\R^d$; the enlarged cube has side less
than $L$, and the same disjoint-volume comparison applies.
\end{proof}

\begin{lemma}[A net with covering radius below twice its spacing]
\label{lem:net}
Let $s\ge2$ be a power of two and, in the torus case, suppose
$L>40s$. An $s$-separated set that $(2s-1)$-dominates can be
computed in
\[
                    O\bigl(ds(5^{6d}+1+\ls n)\bigr)
\]
deterministic rounds.
\end{lemma}

\begin{proof}
Set $S_{-1}=V(G)$. For $j=0,\ldots,\log_2s$ put $r_j=2^j$,
and form an auxiliary graph $H_j$ on $S_{j-1}$ by joining centers
whose maximum-norm distance is at most $r_j$. Let $S_j$ be a
maximal independent set of $H_j$.

When $j=0$ the closed neighborhoods have size at most $3^d$.
For $j\ge1$, the set $S_{j-1}$ is $r_j/2$-separated, so packing
bounds each closed neighborhood by
$(1+2r_j/(r_j/2))^d=5^d$. All these uses on a torus satisfy
$2r_j+r_j/2<L$. Thus each $H_j$ has degree at most $5^d$.
Identifiers give an initial proper coloring with $m=n^C$ labels;
Lemma~\ref{lem:reduction} computes its maximal independent set in
$O(5^{6d}+1+\ls n)$ auxiliary rounds. Auxiliary neighbors can
communicate in $O(dr_j)$ physical rounds. Routes and neighbors
are determined by gathering the corresponding grid neighborhood;
unbounded message sizes allow all routes to be used concurrently.

Maximality ensures that every point of $S_{j-1}$ is within $r_j$
of $S_j$. Induction gives separation $r_j$ and domination radius
at most $\sum_{i=0}^j2^i=2^{j+1}-1$. Finally,
$\sum_{j=0}^{\log_2s}r_j<2s$, proving the time bound and the
claimed properties of the last set.
\end{proof}

The covering radius $2s-1$ is sufficient; we do not need an exactly
maximal $s$-separated set. This is what permits the use of sparse
auxiliary graphs throughout the construction.

\subsection{Choosing boxes with separated faces}\label{sec:boxes}

Fix
\begin{equation}\label{eq:upperparameters}
 s=2^{\lceil\log_2(16d\,25^d)\rceil},\qquad
 D_0=25^d,\qquad R=8s+2.
\end{equation}
In particular $16d25^d\le s<32d25^d$, so $s=2^{O(d)}$.
If $L\le40s$, gather the entire graph in $O(dL)=2^{O(d)}$
rounds and choose a canonical proper three-coloring using the
identifier order. Such a coloring exists: each coordinate path
or cycle has a proper three-coloring, and the sum of its coordinate
colors modulo three is proper on their Cartesian product. A
single-vertex factor causes no difficulty. Exhaustive local search
after gathering is allowed. This branch needs no supplied frame.

Henceforth suppose $L>40s$. Obtain centers $S$ from
Lemma~\ref{lem:net} and form the conflict graph $J$ on $S$, joining
centers at maximum-norm distance at most $R$. Packing gives
\[
 \Delta(J)+1\le(1+2R/s)^d=(17+4/s)^d\le19^d\le D_0.
\]
The torus condition holds because $2R+s=17s+4<L$.
Use Lemma~\ref{lem:reduction} to color $J$ with
$O((D_0+1)^6)$ colors, and process these color classes in increasing
order. Each center $a$ selects an integer radius
\begin{equation}\label{eq:radiusrange}
                         r_a\in\{2s+1,\ldots,4s\}.
\end{equation}
For every previously processed conflict neighbor $b$, every coordinate
$i$, and both signs $\epsilon,\eta\in\{-1,1\}$, require
\begin{equation}\label{eq:faces}
           \dist(a_i+\epsilon r_a,b_i+\eta r_b)\ge2.
\end{equation}
This distance is ordinary coordinate distance in a box and cyclic
distance on a torus. In a box we impose the condition even on face
coordinates outside the input box.

For a fixed $b,i,\epsilon,\eta$, at most three candidate radii
violate~\eqref{eq:faces}: the face coordinates would have difference
$-1,0$, or $1$. On a torus these are three congruences modulo $L$,
each with at most one solution in the candidate interval, whose
length $2s$ is less than $L$. Each previous neighbor thus forbids
at most $12d$ radii. Since
\[
                          2s>12dD_0,
\]
a permissible radius always exists; choose the least one. Centers
processed in the same class are not conflict neighbors and make
their choices independently.

Let $B_a=B_\infty(a,r_a)$ and call
$I_a=B_\infty(a,r_a-1)$ its \emph{interior}. This is the clipped
metric interior just defined, not an extra boundary condition at
the outer boundary of the input box. The interiors cover the grid:
each vertex is within $2s-1$ of a center, whereas $r_a-1\ge2s$.
If the expanded boxes $B_\infty(a,r_a+1)$ and
$B_\infty(b,r_b+1)$ intersect, their centers have distance at most
$8s+2=R$. Their radii therefore satisfy~\eqref{eq:faces}.

\subsection{The boundary-parity conversion}

A vertex $v$ is on the $i$-th border of center $a$ if
\[
        v\in B_a,\qquad v_i=a_i-r_a\ \text{or}\ v_i=a_i+r_a,
\]
with coordinate equality modulo $L$ on a torus. Define
\begin{equation}\label{eq:borderbit}
 b(v)=\#\{(a,i):v\text{ is on the }i\text{-th border of }a\}
                  \pmod2.
\end{equation}
Input-boundary faces created solely by clipping do not contribute
to this count.

\begin{lemma}[Containment of monochromatic components]\label{lem:border}
Every connected component induced by a single value of $b$ is
contained in some $I_a$. It is bipartite and has weak graph
diameter at most $8ds$, meaning that any two of its vertices have
distance at most $8ds$ in the whole graph $G$.
\end{lemma}

\begin{proof}
Consider an edge $uv$ with $u\in I_a$ and $v\notin I_a$.
Only one coordinate, say $j$, changes. Thus $v$ lies on exactly
the $j$-th border of $a$, and $u$ lies on no border of $a$.

No other center $b$ can contribute a $j$-border at $u$ or $v$.
If it did, its expanded box would meet that of $a$, and the
corresponding two parallel faces would have coordinate distance
at most one, contrary to~\eqref{eq:faces}. For a coordinate
$i\ne j$, the border contributions of every other box are
unchanged along $uv$. Indeed the $i$-coordinate is unchanged;
if membership in such a face changed, membership in its box
would change along coordinate $j$, forcing one endpoint onto
a $j$-border of that box. We have just excluded this.
Consequently the unreduced count in~\eqref{eq:borderbit} increases
by exactly one, so $b(u)\ne b(v)$.

Every vertex is in some $I_a$. A path of constant $b$ beginning
there cannot leave $I_a$, because its first exit would cross an
edge of the preceding kind. This proves containment. Coordinate
diameters within $I_a$ are less than $2r_a\le8s$, so ambient
graph distances are at most $8ds$. Finally the containing box
is bipartite. On a torus it lifts to an ordinary lattice box,
since its side length is less than $8s+1<L-1$; no wrapping edge
joins its opposite faces. Clipping by a box boundary also
preserves bipartiteness.
\end{proof}

After computing $b$, each vertex gathers the labeled
radius-$(8ds+1)$ neighborhood. It contains its entire
$b$-component and its boundary, by the weak diameter bound.
Local computation identifies that component, chooses its
minimum-identifier vertex as root, and assigns the parity of
distance to the root \emph{inside the component}. The component
is bipartite, so this gives a well-defined bipartition bit $p(v)$.
Its intrinsic diameter need not be small: the entire component
has already been gathered using ambient graph distance.

Output $(b(v),p(v))\in\{0,1\}^2$. If an edge changes $b$, its
endpoints already differ. Otherwise the endpoints lie in one
component and their bipartition bits differ. The output is a
proper four-coloring.

The conflict-graph color reduction and radius selection cost
\[
                  O\bigl(dR((D_0+1)^6+1+\ls n)\bigr)
\]
physical rounds. Computing the border bits requires gathering
center radii within distance $4s$, at cost $O(ds)$, and the
last component step costs $O(ds)$. Together with the net and
small-side branches, an explicit sufficient bound is
\begin{equation}\label{eq:explicitupper}
 O\bigl(d^2 25^{7d}+d^2 25^d(1+\ls n)\bigr)
                         =2^{O(d)}(1+\ls n).
\end{equation}
This proves the deterministic upper bound with a frame.

\subsection{Recovering the necessary geometry without a frame}
\label{sec:unframed}

The large-side branch has $L>40s\ge5$. In a full box or standard
torus of side at least five, two incident edges belong to different
coordinate axes if and only if they lie together in a four-cycle.
Every such four-cycle is a coordinate square. Thus a radius-two
view identifies the groups of incident edges belonging to one
axis: opposite pairs and, at a box boundary, singletons.

A root may name these groups and signs arbitrarily. Across an
edge, each transverse direction is transported through its unique
coordinate square. The back edge receives the opposite sign,
and the direction of travel is transported by straight continuation,
when the continuation edge exists; a
missing continuation records a boundary. In a neighborhood of
graph radius $t$, radius $t+2$ provides all the squares required
for this transport. To justify consistency, compare the procedure
with the actual Cartesian-product representation of the promised
input: at the root they differ by a signed permutation of axes,
and every square and straight-continuation rule preserves that
agreement. Coordinates are taken modulo the known $L$ on tori.
The accumulated displacement of any path therefore gives the
correct relative coordinates, including when the neighborhood
wraps around the torus.

All geometric tests used above are invariant under signed
permutations: maximum-norm distance, cubes, parallel-face
separation, and the total border count. The chosen centers
depend on auxiliary colors, the radius choice is the least
feasible integer, and the component root depends on identifiers.
Different processors need not use the same axis names.
An additional two layers suffice for each geometric gathering;
the sum of these costs is already covered by
\eqref{eq:explicitupper}. This proves the unframed upper bound.

\section{Removing the size term in the stronger models}\label{sec:transfer}

The geometric construction needs symmetry-breaking labels that
distinguish nearby vertices, rather than globally unique labels.
We now state precisely the published ingredients used to obtain
such labels. These are supporting general-graph results; all
dependence on grid dimension is tracked in the subsequent reduction.

\begin{theorem}[Cited initialization primitives]\label{thm:primitives}
For finite graphs of maximum degree at most a known $D\ge1$:
\begin{enumerate}
\item There is an exact proper $3^D$-coloring distribution with
dependence range and non-signaling locality $O(1+\ls D)$.
The family is consistent under subgraph isomorphisms of enlarged
neighborhoods, in particular under equal named local views on
different inputs.
\item With $n$ and $D$ known, a quantum algorithm starting from a
product state produces a proper $O(D^2)$-coloring in
$O(1+\ls D)$ rounds with global probability at least $1-n^{-a}$
for every fixed requested $a>0$.
\end{enumerate}
The constants in the dependence on $D$ are uniform.
\end{theorem}

Part 1 is the intermediate $3^D$-coloring in the \emph{proof} of
Lemma 6.14 of Akbari et al.~\cite{OnlineQuantum}, obtained from
their rooted-pseudoforest coloring (Lemma 6.11) and composition
lemmas (6.4 and 6.8). We use their corrected February 2026 version.
The statement of their Lemma 6.14 additionally reduces the palette
to $D+1$ at a larger cost; that further step is not used here.
Part 2 is Theorem 6.5 of the September 2026 preprint of
Flin--Li--Suomela~\cite{QuantumBreak}. Its local quantum resources
can depend on $n$ and the accuracy parameter. It does not claim
exact success. These are the two imported symmetry-breaking
theorems in the stronger-model upper bounds.

\begin{lemma}[Power graphs and postprocessing]\label{lem:power}
Let $G^r$ join distinct vertices at $G$-distance at most $r$, where
$r\ge1$ is an integer. A total non-signaling outcome of locality
$t$ on $G^r$ pulls back to locality $O(r(t+1))$ on $G$.
A dependence range $k$ becomes at most $rk$.
Deterministic $b$-round postprocessing adds $b$ to non-signaling
locality and at most $2b$ to dependence range.
A quantum round on $G^r$ can be simulated in $O(r)$ rounds on
$G$, after $O(r)$ setup.
\end{lemma}

\begin{proof}
The radius-$t$ joint view of a set in $G^r$ is determined by
its radius-$r(t+1)$ view in $G$. The extra $r$ determines all
power edges incident to visible vertices, including the
short paths witnessing those edges. Equality of these original
views therefore implies equality of the required power-graph
views and output marginals. If two sets have $G$-distance
greater than $rk$, no path of at most $k$ power edges can join
them; the dependence assertion follows.

A $b$-round postprocessing reads only the $b$-neighborhood of
its output set. The neighborhoods of two sets at distance
greater than $rk+2b$ remain separated by more than $rk$,
which proves the composition bounds. For quantum simulation,
gather the physical radius-$r$ neighborhoods and fix a route
of length at most $r$ for each ordered pair of virtual
neighbors, using identifiers to break ties. Route each virtual
message along its path. All registers may travel concurrently
because capacities are unrestricted. A quantum register is
forwarded along its route, never copied.
\end{proof}

\begin{theorem}[Four colors with no size-dependent time term]
\label{thm:strongupper}
Full balanced boxes and standard tori admit exact four-coloring
outcomes in $\NS$ and $\BD$ with locality $2^{O(d)}$.
They admit $2^{O(d)}$-round quantum four-coloring algorithms
with global probability at least $1-n^{-a}$ for any fixed $a>0$,
without pre-shared entanglement. No frame is required.
\end{theorem}

\begin{proof}
Use $s,R$ from~\eqref{eq:upperparameters} and set
\begin{equation}\label{eq:powerradius}
             r=64ds+2,\qquad D_{\mathrm{pow}}=(2d)^{r+1}.
\end{equation}
The latter is an upper bound on the degree of $G^r$ whenever
$G$ has maximum degree at most $2d$: the number of reachable
vertices is at most $\sum_{j=0}^r(2d)^j<(2d)^{r+1}$.
Every auxiliary edge used by the geometric construction has
physical length at most $dR<r$. Every final monochromatic
component has weak diameter at most $8ds<r$.

For non-signaling and bounded dependence, apply part 1 of
Theorem~\ref{thm:primitives} to $G^r$. Let $X$ denote the resulting
proper coloring, with labels in $[Q]$, $Q=3^{D_{\mathrm{pow}}}$.
By Lemma~\ref{lem:power}, its locality and dependence range
on $G$ are $O(r(1+\ls D_{\mathrm{pow}}))$.

Run the geometric algorithm using $X$ as the initial coloring
in every invocation of Lemma~\ref{lem:reduction}. This is valid
even though the sets of centers are chosen adaptively: all
edges of all auxiliary graphs have length at most $r$, so
restriction of $X$ to any set of retained centers is still a
proper coloring of the relevant auxiliary graph. Its palette
bound remains $Q$; the size term in the algorithm is therefore
$1+\ls Q$ instead of $1+\ls n$. For component bipartitions,
the minimum $X$-label is unique within each component because
its weak diameter is less than $r$. Identifiers may also be
retained for gathering and root choice; their magnitude never
enters the remaining round bound.

If $L\le40s$, the whole grid has diameter at most $dL<r$,
so its $X$-labels are distinct and the whole-graph branch
works within the same bound. Otherwise use the separated-box
construction. Quantitatively,
\[
 r=2^{O(d)},\qquad
 D_{\mathrm{pow}}\le2^{2^{O(d)}},\qquad
 \ls Q=O(1+\ls d).
\]
The initialization and the deterministic postprocessing each
take $2^{O(d)}$ locality. Lemma~\ref{lem:power} proves the same
bound for dependence range. Correctness is exact.

For quantum computation, first handle $n=1$ directly and otherwise
apply part 2 of Theorem~\ref{thm:primitives} once to the entire
power graph, with degree bound
$\min\{D_{\mathrm{pow}},n-1\}$. This gives
$Q=O(D_{\mathrm{pow}}^2)$ initial labels, so again
$\ls Q=O(1+\ls d)$. Simulating the virtual rounds by
Lemma~\ref{lem:power} costs $2^{O(d)}$ physical rounds.
Conditioned on the one event that this initial coloring is
proper, every later step is deterministic and correct.
Thus the original global success guarantee $1-n^{-a}$ is
preserved; no union bound over the later auxiliary computations
is needed.
\end{proof}

For completeness, these constructions have total outcomes on
inputs violating the promise. First delete every edge incident
to a vertex of degree greater than $2d$, a local operation that
changes no legal grid. Form the power graph of the resulting
bounded-degree graph and use the total initialization family.
Run the geometric postprocessing for its fixed prescribed time;
when a local view fails a required geometric test, use a fixed
fallback rule. On all inputs this is a bounded-radius
deterministic postprocessing, whether or not its coloring is
proper. Hence the resulting family satisfies exactly the
cross-input locality convention used in the lower bound.

\section{The classical size term, uniformly in dimension}
\label{sec:size}

For fixed dimension the classical size dependence is familiar.
We include the reduction that keeps its constants independent of
$d$, and a short proof of the deterministic ingredient.

\subsection{A deterministic path and cycle lower bound}

For $m\ge1$, let $S_m(L)$ have as vertices the strictly increasing
$m$-tuples from $[L]$. Join the two consecutive $m$-windows of
each increasing $(m+1)$-tuple. In particular $S_1(L)=K_L$.
If $f$ is a $k$-coloring of $S_m(L)$ for $m\ge2$, assign to
$(x_1,\ldots,x_{m-1})$ the subset
\[
 \{f(x_1,\ldots,x_{m-1},y):y>x_{m-1}\}\subseteq[k].
\]
This is a proper coloring of $S_{m-1}(L)$ with at most $2^k$
colors. Indeed, for an increasing $m$-tuple, the color
$f(x_1,\ldots,x_m)$ belongs to the subset for its prefix
and does not belong to the subset for its suffix: every
extension of the suffix is adjacent to the original tuple.
Empty subsets cause no exception to this argument.
Iterating yields
\begin{equation}\label{eq:shift}
     L\le \underbrace{2^{2^{\cdot^{\cdot^{2^k}}}}}_{m-1
        \text{ exponentiations}}
     \quad\text{whenever }S_m(L)\text{ is }k\text{-colorable}.
\end{equation}

Now fix all public parameters of a deterministic $T$-round
algorithm on an oriented $L$-cycle, with identifiers in $[L]$.
If $L\ge2T+4$, its output on an increasing local ID sequence
defines a function on $S_{2T+1}(L)$. Each pair of adjacent
windows can be embedded in one cycle ID assignment, completing
the unused identifiers arbitrarily. Correctness forces the two
outputs to differ. A $c$-color algorithm therefore induces a
$c$-coloring of this shift graph. By~\eqref{eq:shift},
\[
                         2T\ge\ls L-\ls c.
\]
For an oriented path, place the windows in the interior, leaving
one unused vertex at either end; the same safe condition
$L\ge2T+4$ suffices. If that condition fails, $T$ is already
linear in $L$, which implies the stated iterated-logarithmic
lower bound up to an absolute additive constant. Restricting
to even $L$ does not change the proof. This is the standard
shift-graph argument underlying Linial's lower bound~\cite{Linial}.

\subsection{Private randomness}

We use the following precise published ingredient: for each
fixed $\rho>1/2$, private-randomness three-coloring of directed
$L$-cycles, with distinct adversarial identifiers in $[L]$ and
global success at least $\rho$, requires $\Omega(\ls L)$
rounds~\cite[Theorem 2.1]{Naor}. The proof allows arbitrary local
decision functions for each fixed $L$; knowledge of $L$ and
other fixed public parameters does not weaken the bound.

It extends to any fixed palette $c\ge3$: after a proper
$c$-coloring, process colors $c,c-1,\ldots,4$, replacing each
by an available member of $\{1,2,3\}$. A cycle vertex has at
most two neighbors, and the simultaneously processed vertices
form an independent set. This takes $c-3$ rounds and preserves
success whenever the initial coloring is proper.

Here is an explicit reduction covering paths as well. Run a
putative $T$-round path algorithm on a cycle with the same
public size, interpreting its interior local rule at every
vertex. As in the dimension lower bound, normalize any output
outside $[c]$ to a fixed color. For the analysis only, choose a
uniformly random edge to cut, independently
of all private tapes, and couple the path and cycle runs by
using the same tapes and IDs. If the cycle output has a bad
edge, select one deterministically from the realized vector.
Except for at most $2T+3$ choices of cut edge, that bad edge
and the views producing its outputs are unchanged on the
resulting path. If every path succeeds with probability at
least $2/3$, averaging over cuts gives
\[
 \Prob[\text{cycle failure}]\left(1-\frac{2T+3}{L}\right)
                         \le\frac13.
\]
For $T\le L/16$ and $L\ge24$, the parenthesized factor is at
least $3/4$, so cycle success is at least $5/9>1/2$ and Naor's
bound applies. Otherwise $T$ is already linear in $L$, or $L$
is bounded. Thus paths have the same $\Omega_c(\ls L)$ lower
bound up to an additive constant. The cycle result and its
proof also apply when only even lengths are considered.

\subsection{Simulating a grid by a path or cycle}

One processor of an oriented $L$-cycle can simulate an entire
$(d-1)$-dimensional cross-section of an $L^d$-vertex torus.
Index the $L^{d-1}$ transverse positions by
$j\in\{0,\ldots,L^{d-1}-1\}$. At a cycle processor with ID
$i\in[L]$, give its virtual vertex $j$ the identifier
\[
                           (i-1)L^{d-1}+j+1\in[n].
\]
These are all distinct. Edges in transverse directions are
internal to a simulator; edges in the first direction connect
adjacent simulators. Unbounded messages permit one grid round
to be simulated in one cycle round. The supplied grid frame,
if any, is also simulated consistently. Reading the output at
one fixed transverse position produces a proper $c$-coloring
of the cycle whenever the grid output is proper.

For randomized algorithms, each simulator draws the collection
of private tapes for all its virtual vertices. Collections
at distinct simulators are independent, exactly as required
by the private-randomness lower bound. For boxes use an
oriented $L$-path and transverse boxes. The rule may depend on
$d$, on $n=L^d$, and on the grid type; these are fixed public
parameters in the one-dimensional lower bound. We have proved
\[
                           T\ge a_c\ls L-b_c
\]
with constants independent of $d$.

Finally, under~\eqref{eq:large} we have $d\le\sqrt L$ and
$\log_2L\le\sqrt L$ for $L\ge128$. Thus
$\log_2 n=d\log_2L\le L$ and
\[
                              \ls n\le1+\ls L.
\]
Combining the dimension and size lower bounds, using
$\max\{x,y\}\ge(x+y)/2$, gives
$\Omega_c(d+\ls n)$ up to an additive constant for deterministic
and private-randomness $\LOCAL$. This finishes the lower-bound
claims of Theorem~\ref{thm:main}.

\section{The remaining dimension gap}\label{sec:gap}

For each fixed $c\ge4$, even non-signaling outcomes must see a
distance proportional to $d$, while the construction here uses
a distance exponential in $d$. When $L$ is of order $d^2$,
$\ls n=\ls d+O(1)$, so the lower bound is a genuine dimension
obstruction rather than a disguised size lower bound. In
particular it rules out $O(\log d+\ls n)$ time for a fixed
palette in the classical models and $O(\log d)$ locality
in the stronger models.

There is a concrete limitation of the separated-box proof.
Suppose radius assignment must work for every $s$-separated,
$s$-dominating net, using the interval
$\{2s+1,\ldots,4s\}$ and condition~\eqref{eq:faces} whenever
expanded boxes intersect. Consider the $2^{d-1}$ centers
\[
                         \{0,s+1\}^{d-1}\times\{0\}.
\]
They are $s$-separated and can be extended to a maximal such
set in a sufficiently large finite box, which is then
$s$-dominating. Every pair of their boxes overlaps because
their center distance is $s+1$, while their radii exceed $2s$.
All their positive faces in the last coordinate are at positions
$r_a$. Condition~\eqref{eq:faces} therefore requires these
$2^{d-1}$ radii to be pairwise at least two apart. The interval
of $2s$ consecutive candidate integers accommodates at most
$s$ such radii. Necessarily $s\ge2^{d-1}$.

Thus a radius rule that accommodates arbitrary nets in this
way already needs an exponential scale. This does not prove
an exponential coloring lower bound. A more structured net,
a different covering, or a different conversion of a covering
to a coloring could improve the upper bound. Whether the
dimension dependence for four colors can be polynomial is
not settled by these arguments.

\appendix
\section{An elementary color-reduction subroutine}\label{app:reduction}

We prove Lemma~\ref{lem:reduction} so that every deterministic
step of the geometric upper bound is specified. The following
standard set-system method only needs a polynomial number of
colors in the degree; optimizing its exponent is unnecessary.

Fix integers $D\ge1$ and $m\ge2$. There are sets
$A_1,\ldots,A_m\subseteq[M]$, where
\begin{equation}\label{eq:selectiveM}
                     M\le\left\lceil8(D+1)^2\log(m+1)\right\rceil,
\end{equation}
such that for every $i$ and every
$J\subseteq[m]\setminus\{i\}$ with $|J|\le D$,
\begin{equation}\label{eq:selective}
                         A_i\setminus\bigcup_{j\in J}A_j\ne\varnothing.
\end{equation}
To see this, put each universe element in each set independently
with probability $p=1/(D+1)$. For a fixed pair $(i,J)$, a given
element witnesses~\eqref{eq:selective} with probability
\[
 p(1-p)^{|J|}\ge\frac{1}{D+1}\left(\frac D{D+1}\right)^D
                       \ge\frac1{4(D+1)}.
\]
The last inequality follows from $(1+1/D)^D<e<4$.
There are at most $(m+1)^{D+1}$ pairs $(i,J)$: encode $i$
and list the members of $J$ in $D$ padded positions.
With $M=\lceil8(D+1)^2\log(m+1)\rceil$, a union bound bounds
the probability that any pair has no witness by
\[
 (m+1)^{D+1}\exp\left(-\frac{M}{4(D+1)}\right)
                     \le(m+1)^{-(D+1)}<1.
\]
Hence such a family exists. It is a deterministic, uniformly
computable finite table: enumerate $M=1,2,\ldots$ and all
membership tables, stopping at the first that satisfies the
finite tests~\eqref{eq:selective}. The argument proves that
this search terminates within~\eqref{eq:selectiveM}. Its
computational cost consumes no communication rounds.

Given a proper $[m]$-coloring of a degree-$D$ graph, vertices
exchange their colors once. A vertex of color $i$ chooses
the least element of $A_i$ outside the sets of its neighbors'
colors. This is possible by~\eqref{eq:selective}. If its
neighbor has old color $j$, the chosen element is outside
$A_j$, whereas that neighbor's new color belongs to $A_j$.
Thus the new coloring is proper and uses at most $M$ colors.

Let $b=16(D+1)^2\ge64$. For $m\ge2$, the palette bound is
\[
                              M\le b\log_2m.
\]
Indeed $\log(m+1)\le2\log m$ and $\log2<3/4$, so
$M\le12(D+1)^2\log_2m+1\le16(D+1)^2\log_2m$. If
$m\ge2^{b^2}$, two applications reduce the palette to at most
\[
 b\log_2(b\log_2m)\le\log_2m.
\]
For the inequality, write $x=\log_2m\ge b^2$; the function
$x-b(\log_2b+\log_2x)$ is increasing there and at $x=b^2$
has value $b(b-3\log_2b)\ge0$. Every pair of reductions
therefore removes at least one iterated-logarithm level until
the palette is at most $2^{b^2}$. One further reduction gives
at most $b^3=4096(D+1)^6$ colors. An already smaller palette
can simply be kept. This uses $O(1+\ls m)$ rounds, with an
absolute constant. A proper one-color input is edgeless
and needs no reduction.

Finally process the reduced color classes in order. A vertex
joins an independent set if none of its already processed
neighbors has joined. Simultaneously processed vertices are
nonadjacent. Every excluded vertex has a selected neighbor,
so the resulting independent set is maximal. The sweep takes
at most $4096(D+1)^6$ rounds. This proves the lemma, including
its uniformity and its dependence on the initial palette.

\begin{thebibliography}{99}

\bibitem{Alon}
Noga Alon.
\emph{The chromatic number of random Cayley graphs}.
European Journal of Combinatorics 34(8):1232--1243, 2013.
\url{https://www.cs.tau.ac.il/~nogaa/PDFS/colcayley6.pdf}.

\bibitem{grid}
Sebastian Brandt, Juho Hirvonen, Janne H. Korhonen,
Tuomo Lempi\"ainen, Patric R. J. \"Osterg\aa rd,
Christopher Purcell, Joel Rybicki, Jukka Suomela,
and Przemys\l aw Uzna\'nski.
\emph{LCL problems on grids}.
PODC 2017, pp.\ 101--110.
\url{https://arxiv.org/abs/1702.05456v2}.
See Section 8, Theorem 4 and Lemma 8.

\bibitem{HSW}
Alexander E. Holroyd, Oded Schramm, and David B. Wilson.
\emph{Finitary coloring}.
Annals of Probability 45(5):2867--2898, 2017.
\url{https://arxiv.org/abs/1412.2725}.
See Lemmas 12--14 and Corollary 15.

\bibitem{Linial}
Nathan Linial.
\emph{Locality in distributed graph algorithms}.
SIAM Journal on Computing 21(1):193--201, 1992.
\url{https://doi.org/10.1137/0221015}.

\bibitem{Naor}
Moni Naor.
\emph{A lower bound on probabilistic algorithms for distributive ring coloring}.
SIAM Journal on Discrete Mathematics 4(3):409--412, 1991.
\url{https://doi.org/10.1137/0404036}.
Author manuscript:
\url{https://www.wisdom.weizmann.ac.il/~naor/PAPERS/local_coloring.pdf}.
See Theorem 2.1.

\bibitem{NScolor}
Xavier Coiteux-Roy, Francesco d'Amore, Rishikesh Gajjala,
Fabian Kuhn, Fran\c{c}ois Le Gall, Henrik Lievonen,
Augusto Modanese, Marc-Olivier Renou, Gustav Schmid,
and Jukka Suomela.
\emph{No distributed quantum advantage for approximate graph coloring}.
STOC 2024.
\url{https://arxiv.org/abs/2307.09444v3}.
See Definitions 5.3--5.6 and Theorem 5.7.

\bibitem{OnlineQuantum}
Amirreza Akbari, Xavier Coiteux-Roy, Francesco d'Amore,
Fran\c{c}ois Le Gall, Henrik Lievonen, Darya Melnyk,
Augusto Modanese, Shreyas Pai, Marc-Olivier Renou,
V\'aclav Rozho\v{n}, and Jukka Suomela.
\emph{Online locality meets distributed quantum computing}.
STOC 2025; corrected version, 13 February 2026.
\url{https://arxiv.org/abs/2403.01903v4}.
See Lemmas 6.4, 6.8, 6.11 and the proof of Lemma 6.14.

\bibitem{QuantumBreak}
Maxime Flin, Longcheng Li, and Jukka Suomela.
\emph{Quantum advantage for distributed symmetry breaking}.
Preprint, 22 September 2026.
\url{https://arxiv.org/abs/2609.26788v1}.
See Theorem 6.5.

\end{thebibliography}
\end{document}
